\chapter{Detailed Implementation Guidance for OpenFlight}
\label{app:impl}

This appendix turns the review into engineering guidance at the level of
signal-processing parameters, algorithms, and procedures, citing the
governing sources throughout. Coordinate conventions follow
\cref{ch:params}.

\section{Radar signal chain (OPS243-A)}
\label{app:radardsp}

\subsection{Doppler scaling and resolution budget}

At $f_c = 24.125$\,GHz, \cref{eq:doppler} gives 71.7\,Hz per mph. With
OpenFlight's 30\,kHz I/Q sample rate, the unambiguous span is
$\pm15$\,kHz $\approx \pm209$\,mph --- adequate for ball speeds to
$\sim$200\,mph. Velocity resolution is set by observation time:
a 128-sample window ($4.27$\,ms) gives $\Delta f = 234$\,Hz
$\approx3.3$\,mph per raw bin; zero-padding to 4096 (the current
pipeline) interpolates the peak but does not add information. For
\emph{spin} work the window, not the padding, must grow: resolving
sidebands at $f_{\mathrm{spin}} = S/60$ (e.g.\ 50\,Hz at 3,000\,rpm)
requires windows of $\gtrsim40$\,ms, during which a 150\,mph ball
decelerates and the central line \emph{chirps} --- so the practical
estimator de-chirps first (track the central line per
US8845442's spectral-trace step~\cite{us8845442}, resample the phase to
remove it), then measures the residual modulation.

\subsection{Spin-rate comb estimation}

The patent-documented chain (\cref{sec:radarspin}) maps to this concrete
pipeline~\cite{us8845442,trackmanspinpatentblog}:
\begin{enumerate}\itemsep2pt
\item STFT with 50--75\% overlap over the post-impact 50--250\,ms;
      track the ball ridge (max-SNR bin per frame with continuity
      constraint).
\item De-chirp: mix each frame down by the tracked ridge frequency so
      the ball line sits at DC.
\item Estimate the sideband comb spacing on the de-chirped spectrum via
      cepstrum or harmonic product spectrum --- both are robust to
      missing harmonics; the patent's \emph{qualification} step
      (verify equal spacing across $\ge$3 consecutive frames, and
      consistency of the harmonic-number assignment) is what separates
      real spin from clutter~\cite{us8845442}.
\item Report spin only when the qualified comb persists; otherwise fall
      back to a flagged estimate from club priors
      (\cref{app:priors}) --- the Garmin R10's documented
      behavior~\cite{garminr10accuracy}.
\end{enumerate}
Expected detection physics: sideband amplitude scales with
ball-surface asymmetry~\cite{us8845442}; range balls and clean urethane
balls read weakly (the reason for Titleist RCT metal-tagged
balls~\cite{mygolfspymlm2pro} and FlightScope's metallic-dot
stickers~\cite{mevoteardown}). Logging which ball type produced each
detection will quantify this for OpenFlight's own statistics.

\emph{Patent posture:} this method is claimed by
US8845442 to $\sim$2029 (US10393870 to Dec.\ 2026)~\cite{us8845442};
the FlightScope phase-demodulation alternative (US9868044, dielectric
lens~\cite{us9868044}) is claimed to $\sim$2034. See \cref{sec:fto}.

\subsection{Club-speed extraction}

The pre-impact spectrum contains a velocity \emph{smear} from heel
(slow) to toe (fast), plus shaft returns below. TrackMan resolves this
by reconstructing the head silhouette and reporting the geometric
center~\cite{trackmanclubspeed}; a single-radar approximation that
tracks a stable reference is: gate the last 30--50\,ms before the
trigger timestamp, form the velocity histogram of CFAR-passing bins,
discard the bottom decile (shaft/hosel) and top decile (toe glint), and
report a fixed percentile (median of the remainder). Validate against
the loft-dependent smash ceiling of \cref{sec:smash}
(\cite{tutelmansmash}); reject or flag shots exceeding it, which
independent testing shows is the signature of toe-lock
errors~\cite{leach2017}.

\section{Angle measurement (K-LD7) and the EKF}
\label{app:ekf}

\subsection{Interferometric angles}

The K-LD7's two receive patches implement \cref{eq:interferometry} with
a $\lambda/2$-class baseline; per-bin phase comparison of the two ADC
channels after the range-Doppler FFT yields one angle per detection ---
the same phase-monopulse principle as
US10850179~\cite{us10850179} and the R10's three-receiver
array~\cite{garminr10data}. Two practical constraints from the
commercial art: (i) angle error grows as SNR falls, so weight each
detection by measured SNR; (ii) mechanical alignment dominates the
error budget --- a $1\degs$ mount error is a $1\degs$ bias on every
launch direction, amplified $\sim$1.15$\times$ into face angle
(\cref{eq:faceinversion}).

\subsection{Recommended filter}

Implement the \cref{sec:ekf} smoother concretely as:
state $\vect{x} = (\vect{p}, \vect{v})$ (add spin states later);
process model \cref{eq:eom} with Smits--Smith
coefficients~\cite{smitssmith}; measurements: OPS243-A radial speed
($\dot r$, high rate, low noise), K-LD7 angle+range detections (lower
rate, SNR-weighted covariance). Run a forward EKF then an RTS
(Rauch--Tung--Striebel) backward smoother over the 50--150\,ms burst,
and evaluate the smoothed state at the sound-trigger timestamp minus
the acoustic delay --- this back-extrapolation is how commercial radars
report launch conditions robustly despite early-flight
clutter~\cite{us8845442,us10338209}. The innovation-gating step of the
EKF doubles as the outlier rejector for multipath and club returns.
The same filtering architecture extends from point tracking (the ball)
to rigid-body tracking (the club) via the screw-theoretic formulation
of \cref{app:screw}, which is the recommended estimation layer for the
IWR6843 integration.

\subsection{Alignment calibration procedure}

Adopt the commercial patterns: (i) Uneekor-style floor chart --- a
printed target-line chart at known positions establishes the
target-line azimuth for both K-LD7s~\cite{uneekorcalib}; (ii)
FlightScope-style known-trajectory check --- roll or swing balls along
a surveyed line and verify reported directions
(US10338209 uses a Doppler simulator for the same
purpose~\cite{us10338209}); (iii) record the mount pose in the session
log so data from different setups is never silently mixed.

\section{Derived club data, with priors}
\label{app:priors}

Implement the D-plane module once, use it both ways
(\cref{sec:dplane,sec:faceangle}):
\begin{itemize}\itemsep2pt
\item \textbf{Forward} (simulation/validation): Tutelman's calibrated
closed form, \cref{eq:obliqueness,eq:tutelmanlaunch}, plus spin from
\cref{eq:spinrate}~\cite{tutelman3d}.
\item \textbf{Inverse} (reporting): face angle from
\cref{eq:faceinversion} with $w_f$ interpolated in dynamic loft between
0.87 (driver) and 0.75 (wedge)~\cite{trackmanballflightlaws,tutelman3d};
dynamic loft analogously from launch angle and attack angle. Tag both
as \emph{derived} per \cref{tab:hierarchy}.
\item \textbf{Priors table} per club type (driver\ldots{}wedge): static
loft, typical dynamic-loft delta, spin-loft range, smash ceiling
(\cref{sec:smash}), PGA/LPGA reference deliveries~\cite{trackman40params}.
Used for: spin fallback estimates, outlier gating, and the smash sanity
check.
\item \textbf{Gear-effect bound}: when impact location is unknown,
attach an uncertainty of up to $\pm6\degs$ of spin-axis tilt per
0.14\,in of possible driver miss~\cite{perfectgolfswing,tutelmangear} to
any derived face-to-path interpretation --- and surface it in the UI
rather than hiding it.
\end{itemize}

\section{Optical module (Phase 2) design parameters}
\label{app:optical}

The expired Wintriss patents~\cite{us7292711} plus PiTrac's published
design~\cite{pitrac} fix the working parameters:
\begin{itemize}\itemsep2pt
\item \textbf{Sensor}: Raspberry Pi Global Shutter camera (IMX296,
$1456\times1088$); global shutter is non-negotiable (rolling shutter
skews a 150\,mph ball by several pixels per row-time).
\item \textbf{Strobed multi-exposure}: $N=3$--5 IR pulses
(850\,nm, tens of \si{\micro\second} each) inside one long exposure
freeze $N$ ball images per frame; at 150\,mph the ball moves 6.7\,cm/ms,
so pulse spacing of $\sim$300--500\,\si{\micro\second} spaces images
2--3\,cm apart in a 30\,cm capture volume --- matching the commercial
capture geometry~\cite{foresightgcquad,pitrac}.
\item \textbf{Trigger}: reuse the SEN-14262 sound trigger; its
$\sim$10\,\si{\micro\second} latency is far inside the strobe-timing
budget, solving the triggering problem the Wintriss patent addressed
with a microphone+radar pair~\cite{us7292711}.
\item \textbf{Ball detection}: Hough circles on the strobed frame
(PiTrac-proven~\cite{pitrac}); depth from the known 42.67\,mm diameter
(the Wintriss mono-camera range cue~\cite{us7292711}) or from a second
camera via \cref{eq:depth}.
\item \textbf{Spin}: register the dimple texture between successive
ball images over $SO(3)$ (\cref{eq:rotangle}); Gabor-filter
pre-enhancement of dimples is documented in
US12401909~\cite{us12401909}; with $\Delta t \approx 400$\,\si{\micro
\second}, 3,000\,rpm is only $7.2\degs$ of rotation --- comfortably
below the aliasing limit, and small enough that a local search around
the D-plane-predicted axis converges quickly
(\cref{sec:dimplespin}). This is the public-domain
(post-2025) Wintriss/Foresight method~\cite{us7292711,foresightspherical}.
\item \textbf{Impact location / measured face angle} (later): either
face fiducials per the expired Acushnet recipes
(US5501463/US6758759: dots + stereo pose)~\cite{us5501463,us6758759},
or an overhead second camera per the Uneekor
geometry~\cite{uneekoreyexo}.
\end{itemize}

\section{Flight model and validation protocol}

Ship Smits--Smith (\cref{eq:smits}) with spin
decay~\cite{smitssmith}, version the coefficients, and calibrate
against measured trajectories using the US6186002 fitting
approach~\cite{us6186002} --- outdoor sessions with the MLM2PRO (or
simple carry ground-truth) provide the data. Validation against the
MLM2PRO should follow the Leach protocol
structure~\cite{leach2017}: per-club shot blocks, Bland--Altman limits
of agreement per parameter, ball data compared directly, spin compared
only on RPT/RCT balls (where the MLM2PRO's spin is
camera-measured~\cite{mygolfspymlm2pro}), and club speed compared with
an explicit reference-point caveat (\cref{sec:radarclub}). Log raw
I/Q and (later) raw frames for every shot so algorithm changes can be
replayed against history --- the practice that made this review's
accuracy analysis possible for the commercial units is exactly what an
open project can do better.
